\clearpage
\section{Application coherence and the remaining copying case}
\label{sec:copying-application-coherence}
Application now has complete successor-state coherence in the original combined
copying machine. Two actual rewrites at the same application request produce
recursively isomorphic complete graphs and corresponding resumed tokens.
The proof includes arbitrary renaming, plugging and simultaneous expansion,
and permits the two matches to describe the actual boxed body by different
recursively isomorphic graphs. Together with the preceding store and arithmetic
results, this milestone isolated copying as the sole remaining family for
successor coherence. Section~\ref{sec:copying-successor-coherence} now resolves
that case and proves complete successor determinism up to representation.
Refocusing, specimen locality and the template argument remain open.

The operational rule and its stability premise remain unchanged. Opening a
lambda exposes its full padded body, preserves its ordered variable and atom
attachments, and connects the first bound variable to the stable argument.
It does not simplify the Figure~16 distributors. In particular, an original
translated variable requests copying after beta; it does not immediately
become its argument. The two terms in Proposition~12.2 are unchanged.

The quantitative result applies to every successful run of the full machine:
application contributes exactly six transitions before the actual opened-body
state. The same law holds under an arbitrary serial host. Constant and
nested-lambda bodies have exactly seven successful transitions, while a
variable body has seven plus the successful cost of its remaining copy request.
No body-termination assumption occurs in the six-step law.

The proof dependency chain is: unique source frame and box body; an explicit
presentation of all surviving and newly exposed occurrences; a realization of
that presentation by the original beta reduct; transport through every rewrite
constructor; comparison of independently represented bodies; whole-state
coherence; exact costs; and the exhaustive copying-only classification.
The presentations and certificates below are mathematical proof data derived
from actual rewrites. They add no premise to the transition relation.

Appendix~B.1 removes application from its outstanding successor comparisons.
Its four broad work packages remain open. The current assessment after the
subsequent copying milestone is in Section~\ref{sec:copying-successor-coherence};
the dependent template and final observational arguments are still unresolved.

\Needspace{12\baselineskip}
\subsection{Presenting the complete result}

\termdef{opening-presentation}
A \emph{partial incidence presentation} $P$ consists of ambient sets $V_P$ and
$E_P$, vertex and edge labels, predicates selecting the retained vertices and
edges, and ordered source and target lists in the ambient port space.
The ambient data need not itself be a well-formed graph: incidence lists for
deleted occurrences are harmless proof data. This allows one space to name
both the old source occurrences and all newly exposed body occurrences.

\lean{Copying/OpeningPresentation}{10}{18}{Opening.Presentation}

\termdef{opening-realization}
A \emph{realization} of $P$ by a genuine hypernet $H$ consists of injections
$j_V:V_H\to V_P$ and $j_E:E_H\to E_P$ whose images are exactly the selected
occurrences. It preserves internal wire labels, relates edge labels by the
recursive relation \texttt{LabelEquiv}, and satisfies
\[
\widehat j_V(\operatorname{src}_H(e))=\operatorname{src}_P(j_E(e)),\qquad
\widehat j_V(\operatorname{tgt}_H(e))=\operatorname{tgt}_P(j_E(e)).
\]
Port extension fixes the external input and output indices, and acts by $j_V$
on internal vertices. The equations are equalities of ordered lists.

\lean{Copying/OpeningPresentation}{21}{31}{Opening.Realization}

Two realizations of the same presentation determine bijections between their
vertex sets and between their edge sets. Send an occurrence to the unique
occurrence with the same ambient image; exact image equality gives existence,
and injectivity gives uniqueness and the inverse equations. The vertex and
edge commuting lemmas express equality of these images. The port commuting
lemma follows by cases on internal, input and output ports.

\lean{Copying/OpeningPresentation}{36}{38}{Opening.Realization.vertexBijection}

\lean{Copying/OpeningPresentation}{39}{41}{Opening.Realization.edgeBijection}

\lean{Copying/OpeningPresentation}{43}{45}{Opening.Realization.vertex_commute}

\lean{Copying/OpeningPresentation}{47}{49}{Opening.Realization.edge_commute}

\lean{Copying/OpeningPresentation}{51}{53}{Opening.Realization.port_commute}

These bijections constitute a recursive graph isomorphism. Labels agree
through their common ambient labels; recursive edge-label comparisons compose
by symmetry and transitivity. To compare source or target lists, embed both
lists in the presentation, use the realization equations and commuting laws,
and cancel the injective port map. If the resumed tokens have the same ambient
image, the vertex bijection also maps one token to the other; with both modes
set to search this gives a state isomorphism.

\lean{Copying/OpeningPresentation}{56}{75}{Opening.Realization.iso}

\lean{Copying/OpeningPresentation}{77}{82}{Opening.Realization.stateIso}

A realization can be pulled back along any recursive isomorphism $K\cong H$.
Compose its injections with the isomorphism's bijections. The inverse maps
supply all required preimages, and composing the label and incidence laws
proves that $K$ realizes the same presentation.

\lean{Copying/OpeningPresentation}{85}{101}{Opening.Realization.rename}

\Needspace{12\baselineskip}
\subsection{The source frame and its actual body}

\termdef{application-source-match}
An \emph{application source match} records three source edges: the application
$a$, its lambda $\ell$, and the lambda's box $b$. It also records internal
function and thunk wires $f,t$, and the argument attachment $q$. The application
head is at the requested token $v$, and the ordered incidences are
\[
\operatorname{tgt}(a)=[f,q],\quad
\operatorname{src}(\ell)=[f],\quad \operatorname{tgt}(\ell)=[t],\quad
\operatorname{src}(b)=[t].
\]
The lambda label is the actual primitive lambda. The box label is recursively
related to the actual body $B$, with one bound term variable, $k$ free term
variables and $h$ free atoms. Equality casts reconcile $k+1$ with $1+k$; they
do not permute the ordered ports. Stability is required by the actual local
transition, not duplicated in this auxiliary source-shape record.

\lean{Copying/ApplicationMatch}{9}{23}{Application.Match}

Two matches of the same source and requested port have identical frame data.
Unique head selection identifies their application edges. Comparing the two
ordered target lists identifies the function wire and argument port. Unique
source ownership identifies the lambda edge at that function wire; its target
identifies the thunk wire, and source ownership again identifies the box edge.
The actual label of this unique box then relates the two candidate bodies.
The relation \texttt{BoxRelated} fixes their binder/interface metadata and
supplies a recursive isomorphism of the full bodies.

\lean{Copying/ApplicationMatch}{30}{46}{Application.Match.shape_unique}

\lean{Copying/ApplicationMatch}{49}{56}{Application.Match.bodies_related}

The box signature proves that the list $q::\operatorname{tgt}(b)$ has length
$1+|\Gamma|+|\Delta|$. This is exactly the opened body's output interface.
The body-port map sends a body internal vertex to a new occurrence, its unique
input to the old request port $v$, and output $i$ to entry $i$ of this list.
Thus output zero reaches the argument; the remaining outputs retain the
original box's ordered free term and atom attachments.

\lean{Copying/ApplicationMatch}{58}{64}{Application.Match.output_length}

\lean{Copying/ApplicationMatch}{67}{70}{Application.Match.bodyPort}

The presentation for a match has ambient sets
\[
V_G+V_{\operatorname{open}(B)},\qquad E_G+E_{\operatorname{open}(B)}.
\]
It retains every old vertex except $f,t$, every old edge except $a,\ell,b$,
and every newly exposed body vertex and edge. Old occurrences retain their
labels and ordered incidence lists; body occurrences use their actual labels
and the body-port map just defined. Here $\operatorname{open}(B)$ is precisely
the existing depth-raised, interface-cast body used by the original reduct.

\lean{Copying/ApplicationMatch}{73}{89}{Application.Match.presentation}

\Needspace{12\baselineskip}
\subsection{The original beta reduct realizes the presentation}

The local source match selects the application, lambda and box in the
existing \texttt{before} graph. Its fields reduce directly to the source
construction. The local realization maps the reduct's argument-connection
wire to the original argument wire, its body internals to new occurrences,
and its argument internals to the corresponding old source occurrences.
The edge map similarly sends body edges to new occurrences and argument edges
to old occurrences.

\lean{Copying/ApplicationOpeningLocal}{8}{22}{Application.localMatch}

\lean{Copying/ApplicationOpeningLocal}{25}{33}{Application.openingVertices}

\lean{Copying/ApplicationOpeningLocal}{35}{41}{Application.openingEdges}

The vertex and edge maps are injective by disjointness of the tagged summands
and injectivity within each summand. Their exact-image lemmas are stronger
than injectivity: every old occurrence except the two frame wires and three
frame edges has a preimage, and every body occurrence has a preimage.
The proofs enumerate the finite frame indices and the host/body summands;
no assumption about the body's size, termination or absence of nested boxes
is used.

\lean{Copying/ApplicationOpeningLocal}{43}{57}{Application.opening_vertices_injective}

\lean{Copying/ApplicationOpeningLocal}{59}{69}{Application.opening_edges_injective}

\lean{Copying/ApplicationOpeningLocal}{71}{101}{Application.opening_vertex_range}

\lean{Copying/ApplicationOpeningLocal}{103}{127}{Application.opening_edge_range}

The serial-output lemma states that an output port of a right serial
component remains the corresponding output of the serial composite. Applying
it through the local construction proves that the actual box's target list is
the free environment ports, in order. The body-port comparison then checks
internal, input and output cases; at output zero it obtains the argument
connection, and at successor indices it obtains the free environment port.
The argument-port comparison follows the same cases using the existing
argument embedding and right-serial output law.

\lean{Copying/ApplicationOpeningLocal}{131}{137}{Application.right_serial_output}

\lean{Copying/ApplicationOpeningLocal}{140}{152}{Application.local_box_targets}

\lean{Copying/ApplicationOpeningLocal}{155}{168}{Application.opening_body_port}

\lean{Copying/ApplicationOpeningLocal}{171}{188}{Application.opening_argument_port}

The complete local realization combines the injections and exact images.
For a body edge, the source and target equations follow by mapping the
body-port comparison over its ordered lists. For an argument edge, use the
argument-port comparison instead. Labels are exactly the construction's
labels, including every nested body box. This proves that the original
\texttt{after} graph realizes the full source-relative presentation.

\lean{Copying/ApplicationOpeningLocal}{191}{235}{Application.opening_realization}

\Needspace{12\baselineskip}
\subsection{Occurrences and changes of presentation}

An occurrence map consists of maps on ports, internal vertices and edges.
The latter two maps are injective, internal ports map to their indicated
internal occurrences, wire labels are preserved, edge labels are recursively
related, and source and target lists commute with the port map. It does not
require arbitrary boundary attachments to be injective.

\lean{Copying/OpeningOccurrence}{8}{18}{Opening.Occurrence}

Forgetting the additional recursive-label and internal-vertex information
gives the existing primitive matching embedding. A recursive isomorphism gives
an occurrence map by its bijections. Plugging gives one by the selected
component summand; simultaneous expansion gives one by tagging each component
occurrence with the selected edge index. Their label and incidence laws are
those of the actual graph constructions.

\lean{Copying/OpeningOccurrence}{22}{28}{Opening.Occurrence.matchEmbedding}

\lean{Copying/OpeningOccurrence}{30}{40}{Opening.Occurrence.ofIso}

\lean{Copying/OpeningOccurrence}{42}{54}{Opening.Occurrence.plug}

\lean{Copying/OpeningOccurrence}{56}{68}{Opening.Occurrence.expand}

Mapping a source match along an occurrence map preserves all frame equations.
Primitive head and lambda labels use the matching embedding, while the box
comparison uses the full recursive label relation. The ambient origin-port
map sends old occurrences through the occurrence map and keeps new body
occurrences fixed; external ports follow the occurrence's boundary attachment.
It commutes with embedding an old port. The body-port compatibility lemma
holds internally and at the body input by definition, and at each output by
mapping the ordered argument-plus-box-target list.

\lean{Copying/OpeningOccurrence}{78}{91}{Application.Match.map}

\lean{Copying/OpeningOccurrence}{94}{100}{Application.Match.originPort}

\lean{Copying/OpeningOccurrence}{102}{104}{Application.Match.origin_old}

\lean{Copying/OpeningOccurrence}{106}{118}{Application.Match.bodyPort_map}

An isomorphism of presentations consists of ambient vertex and edge
bijections preserving the retention predicates, labels and ordered incidence
lists. Composing a realization with these bijections transports it to the
new presentation. Inverses prove the exact-image conditions; the remaining
conditions follow by composition.

\lean{Copying/OpeningPresentationIso}{8}{16}{Opening.Presentation.Iso}

\lean{Copying/OpeningPresentationIso}{19}{40}{Opening.Realization.transport}

For a recursively renamed source, the origin-port map is the port extension
of the sum of its vertex bijection with the identity on new body vertices.
This gives the source-renaming presentation isomorphism. Old label and
incidence equations come from the source isomorphism; new body incidences use
the body-port compatibility lemma.

\lean{Copying/OpeningPresentationIso}{48}{53}{Application.Match.origin_iso}

\lean{Copying/OpeningPresentationIso}{56}{97}{Application.Match.renamePresentation}

For two matches of one source, a recursive isomorphism between their opened
bodies also induces a presentation isomorphism. It is the identity on old
occurrences and the body isomorphism on new ones. The body-port comparison
uses unique frame and argument data: internal ports follow the body bijection,
inputs agree at $v$, and outputs use the identical ordered attachment list.
Retention predicates agree because the removed frame occurrences are unique.
Labels and ordered incidences then follow from the body isomorphism.

\lean{Copying/OpeningPresentationIso}{100}{110}{Application.Match.bodyPort_congr}

\lean{Copying/OpeningPresentationIso}{113}{146}{Application.Match.bodyPresentation}

\Needspace{12\baselineskip}
\subsection{Complete certificates and independent bodies}

\termdef{application-opening-certificate}
An \emph{opening certificate} for $G,v\longrightarrow H,w$ records a source
match, a realization of its presentation by $H$, and the token equation
$\widehat j_V(w)=\operatorname{inl}(v)$. Thus it determines the full graph and
where execution resumes. The original local rule supplies such a certificate
using the local realization and its input-token equation.

\lean{Copying/ApplicationOpeningCertificate}{8}{12}{Application.OpeningCertificate}

\lean{Copying/ApplicationOpeningCertificate}{15}{19}{Application.local_certificate}

Equal interface casts can be removed from a recursive graph isomorphism:
eliminate the interface equalities, after which the original isomorphism has
the required type. To rename both sides of a certificate, transport its
presentation along the source isomorphism and pull back its realization along
the inverse successor isomorphism. Port-map composition and the successor
inverse equation prove the renamed token equation.

\lean{Copying/ApplicationOpeningCertificate}{22}{24}{Application.cast_iso_reflect}

\lean{Copying/ApplicationOpeningCertificate}{31}{42}{Application.OpeningCertificate.rename}

Two opening certificates for the same source request are coherent even when
their body graphs and interface metadata are supplied independently. The
unique source box first identifies the metadata and relates the cast bodies.
Remove the casts, raise the resulting recursive isomorphism through depth
padding, and restore the opened-body interface cast. This gives the required
isomorphism between the actual opened bodies. Transport one realization
through its induced presentation isomorphism, then compare the two
realizations of one presentation. Their token equations agree on the old
source port, yielding the complete state isomorphism.

\lean{Copying/ApplicationOpeningCertificate}{45}{59}{Application.OpeningCertificate.coherent}

\Needspace{12\baselineskip}
\subsection{Transport through arbitrary plugging}

The auxiliary shift map sends an old component occurrence to the old
component summand of the host, and keeps a new body occurrence new:
$A+B\to(C+A)+B$. It is injective. Its old- and new-image equivalences identify
preimages in the corresponding summands, and its disjointness lemma excludes
old host occurrences from its image. All four proofs split the sum constructors.

\lean{Copying/ApplicationOpeningPlug}{8}{10}{Opening.shift}

\lean{Copying/ApplicationOpeningPlug}{12}{13}{Opening.shift_injective}

\lean{Copying/ApplicationOpeningPlug}{15}{16}{Opening.shift_old}

\lean{Copying/ApplicationOpeningPlug}{18}{19}{Opening.shift_new}

\lean{Copying/ApplicationOpeningPlug}{21}{21}{Opening.shift_not_host}

The plugged vertex and edge injections keep old host occurrences in their
host summands and shift the selected component's certified images. The
injectivity lemmas use host/component disjointness, shift injectivity and the
original certificate's injectivity. The range lemmas split host, old component
and new body occurrences: the host is retained in full, and the component and
body cases reduce to the original exact-image equations.

\lean{Copying/ApplicationOpeningPlug}{32}{34}{Application.OpeningCertificate.plugVertices}

\lean{Copying/ApplicationOpeningPlug}{36}{38}{Application.OpeningCertificate.plugEdges}

\lean{Copying/ApplicationOpeningPlug}{40}{45}{Application.OpeningCertificate.plug_vertices_injective}

\lean{Copying/ApplicationOpeningPlug}{47}{52}{Application.OpeningCertificate.plug_edges_injective}

\lean{Copying/ApplicationOpeningPlug}{54}{61}{Application.OpeningCertificate.plug_vertex_range}

\lean{Copying/ApplicationOpeningPlug}{63}{70}{Application.OpeningCertificate.plug_edge_range}

The host-port lemma preserves the same host attachment under the new
injection. The component-port lemma commutes with insertion and the ambient
origin-port map. Both follow by port cases: internal component ports shift,
while boundary ports use the host-port law. The source and target presentation
lemmas treat old edges using the occurrence's list equations and new edges
using the body-port compatibility law.

\lean{Copying/ApplicationOpeningPlug}{72}{74}{Application.OpeningCertificate.plug_host_port}

\lean{Copying/ApplicationOpeningPlug}{76}{89}{Application.OpeningCertificate.plug_component_port}

\lean{Copying/ApplicationOpeningPlug}{91}{102}{Application.OpeningCertificate.plug_source}

\lean{Copying/ApplicationOpeningPlug}{104}{115}{Application.OpeningCertificate.plug_target}

The plugged realization uses these maps and range laws. Host edge labels and
incidences are retained; component labels come from the original realization,
and their incidences follow by composing the component-port and presentation
list equations. The resulting opening certificate uses the mapped source
match; its token equation is the component-port equation applied to the
original token equation. Thus arbitrary typed plugging preserves complete
opening certificates.

\lean{Copying/ApplicationOpeningPlug}{118}{162}{Application.OpeningCertificate.plugRealization}

\lean{Copying/ApplicationOpeningPlug}{164}{171}{Application.OpeningCertificate.plug}

\Needspace{12\baselineskip}
\subsection{Transport through simultaneous expansion}

The expansion proof first handles dependent occurrence sets. Fix component
sets $F_e,K_e$, a selected index $h$, a map $K_h\to F_h+N$, and bijections
$K_e\cong F_e$ for $e\ne h$. The tag map sends an old selected occurrence to
its pair $(h,x)$ and leaves a new occurrence in $N$. It is injective by
disjointness and injectivity of dependent-pair construction.

\lean{Copying/OpeningFiber}{11}{13}{Opening.Fiber.tag}

\lean{Copying/OpeningFiber}{15}{20}{Opening.Fiber.tag_injective}

The lifted map on $\sum_e K_e$ applies the selected map and tagging at $h$,
and the given bijection at every other index. Its same-index and other-index
lemmas expose these two clauses. A selected tag cannot be an old occurrence
at another index: equality of dependent pairs would equate their indices.
The lifted map is injective by four cases on whether the two indices equal
$h$. The two mixed cases use this disjointness; the selected case uses the
selected injection; the other case uses index equality followed by the
corresponding bijection's injectivity.

\lean{Copying/OpeningFiber}{22}{25}{Opening.Fiber.lift}

\lean{Copying/OpeningFiber}{27}{28}{Opening.Fiber.lift_same}

\lean{Copying/OpeningFiber}{29}{30}{Opening.Fiber.lift_other}

\lean{Copying/OpeningFiber}{32}{35}{Opening.Fiber.tag_not_other}

\lean{Copying/OpeningFiber}{37}{57}{Opening.Fiber.lift_injective}

The range at an old selected occurrence is exactly the range of the selected
map at that old occurrence. The same is true for a new occurrence. In each
forward direction, an index other than $h$ is ruled out by its tag; the
reverse direction supplies index $h$. Every old non-selected occurrence has
a preimage obtained using its component bijection's inverse.

\lean{Copying/OpeningFiber}{59}{70}{Opening.Fiber.range_same}

\lean{Copying/OpeningFiber}{72}{82}{Opening.Fiber.range_extra}

\lean{Copying/OpeningFiber}{84}{86}{Opening.Fiber.range_other}

For actual expansion, the unchanged-family premise supplies a recursive
isomorphism $K_e\cong F_e$ at each non-selected index. Apply the dependent
construction separately to vertices and edges, using the opening realization
at the selected index. Whole-graph vertices retain the outer host summand and
shift the component fibers; whole-graph edges are just the edge fibers.
Their injectivity follows from the dependent injection lemma and host/fiber
disjointness.

\lean{Copying/ApplicationOpeningExpansion}{14}{15}{Application.OpeningCertificate.otherIso}

\lean{Copying/ApplicationOpeningExpansion}{17}{18}{Application.OpeningCertificate.fiberVertices}

\lean{Copying/ApplicationOpeningExpansion}{19}{20}{Application.OpeningCertificate.fiberEdges}

\lean{Copying/ApplicationOpeningExpansion}{21}{23}{Application.OpeningCertificate.expansionVertices}

\lean{Copying/ApplicationOpeningExpansion}{24}{24}{Application.OpeningCertificate.expansionEdges}

\lean{Copying/ApplicationOpeningExpansion}{26}{31}{Application.OpeningCertificate.expansion_vertices_injective}

\lean{Copying/ApplicationOpeningExpansion}{32}{34}{Application.OpeningCertificate.expansion_edges_injective}

The expanded range lemmas keep every outer host vertex and every occurrence
of a non-selected component. At the selected component they reduce to its
certificate's exact images; new body occurrences use the new-range lemma.
Distinct component indices ensure that none of the unchanged occurrences is
mistaken for a deleted application, lambda, box, function wire or thunk wire.

\lean{Copying/ApplicationOpeningExpansion}{36}{55}{Application.OpeningCertificate.expansion_vertex_range}

\lean{Copying/ApplicationOpeningExpansion}{57}{74}{Application.OpeningCertificate.expansion_edge_range}

The host-port equation fixes host attachments. The selected-component port
equation uses the certificate and its ambient origin map; the other-component
port equation uses the corresponding recursive isomorphism. Internal cases
use the appropriate dependent-map clause, and boundary cases reduce to the
host-port equation. The expanded presentation's source and target laws then
follow for old selected edges from the occurrence equations and for new body
edges from body-port compatibility.

\lean{Copying/ApplicationOpeningExpansion}{76}{78}{Application.OpeningCertificate.expansion_host_port}

\lean{Copying/ApplicationOpeningExpansion}{80}{90}{Application.OpeningCertificate.expansion_selected_port}

\lean{Copying/ApplicationOpeningExpansion}{92}{100}{Application.OpeningCertificate.expansion_other_port}

\lean{Copying/ApplicationOpeningExpansion}{102}{113}{Application.OpeningCertificate.expansion_source}

\lean{Copying/ApplicationOpeningExpansion}{115}{126}{Application.OpeningCertificate.expansion_target}

The expanded realization combines these injections, images and port laws.
At a selected edge it uses the original realization's recursive label and
incidence equations. At every other edge it uses the unchanged component's
recursive isomorphism. Its token equation follows from the selected port law
and the original certificate. This constructs a certificate for every actual
simultaneous-expansion rewrite, preserving all other component contents.

\lean{Copying/ApplicationOpeningExpansion}{128}{192}{Application.OpeningCertificate.expansionRealization}

\lean{Copying/ApplicationOpeningExpansion}{194}{201}{Application.OpeningCertificate.expand}

\Needspace{12\baselineskip}
\subsection{Derivation from the unchanged transition relation}

An opening witness hides the actual opened box's depth, interface counts and
body graph while retaining its full certificate. Two witnesses for one source
request are coherent by unpacking their certificates and applying the
independent-body comparison. This existential packaging permits induction over
the machine's depth-indexed rewrite relation.

\lean{Copying/ApplicationCoherence}{9}{12}{Application.OpeningWitness}

\lean{Copying/ApplicationCoherence}{14}{19}{Application.OpeningWitness.coherent}

For an actual local transition recognized as application, the beta case
supplies the original local certificate. Every other local constructor is
excluded by unique head recognition: its primitive label differs from
application. Stability is still present in the beta constructor.
Induction over an arbitrary actual rewrite then produces its witness. The
local case is just proved. For renaming, plugging and expansion, recognition
reflects into the selected component, the induction hypothesis produces its
witness, and the corresponding certificate transport constructs the result.

\lean{Copying/ApplicationCoherence}{26}{55}{Combined.local_application_certificate}

\lean{Copying/ApplicationCoherence}{59}{74}{Combined.rewrite_application_certificate}

\termdef{application-successor-coherence}
\emph{Application successor coherence} states: if a source request is at an
application head and has two actual rewrite successors, those complete search
states are recursively isomorphic. Derive a witness from each rewrite and
compare them. For full machine steps, a search transition from rewrite mode is
impossible, so both steps reduce to the rewrite theorem. No global determinism
assumption, template law or body-termination assumption is used.

\lean{Copying/ApplicationCoherence}{77}{81}{Combined.application_rewrite_deterministic}

\lean{Copying/ApplicationCoherence}{83}{90}{Combined.application_step_deterministic}

For a chosen actual application step $s\to t$, the successful-cost spectrum
satisfies $T_n(s)\iff\exists m\,(n=m+1\land T_m(t))$. A successful zero-step
run cannot start at a state with a transition. Any nonempty successful run
has a first successor isomorphic to $t$, and execution-isomorphism invariance
transports its remaining cost; the reverse direction prepends the chosen step.
Unbounded application steps lift finite-depth steps, so their successors are
related by the existing representation relation. Representation equivalence
then gives their machine equivalence for every cost order.

\lean{Copying/ApplicationCoherence}{92}{95}{Combined.application_cost_spectrum}

\lean{Copying/ApplicationCoherence}{97}{104}{Combined.application_unbounded_deterministic}

\lean{Copying/ApplicationCoherence}{106}{109}{Combined.application_unbounded_equiv}

\Needspace{12\baselineskip}
\subsection{The historical reduction to the copying case}

\termdef{copying-coherence-frontier}
The \emph{copying coherence frontier} is an exhaustive theorem about the
existing machine. Two actual rewrites either have recursively isomorphic
successors or share a term-contraction head. Start from the earlier exhaustive
classification into coherent, application or copying cases, and discharge the
application case by the new theorem. For full steps, search successors are
unique, rewrite/search overlap is impossible, and the rewrite classification
applies otherwise.

\lean{Copying/CombinedCopyingFrontier}{9}{16}{Combined.rewrite_coherent_or_copying}

\lean{Copying/CombinedCopyingFrontier}{18}{28}{Combined.step_coherent_or_copying}

Consequently, excluding a rewrite-mode term-contraction request suffices for
successor coherence and for the exact one-step successful-cost decomposition.
The two corollaries resolve the copying alternative and apply the existing
coherent-step cost lemma. The unbounded classification follows by unpacking
both lifted steps and lifting the resulting state isomorphism. None of these
results alone resolves the copying alternative; Section~\ref{sec:copying-successor-coherence} supplies its complete proof.

\lean{Copying/CombinedCopyingFrontier}{30}{33}{Combined.step_deterministic_away_copying}

\lean{Copying/CombinedCopyingFrontier}{35}{39}{Combined.cost_spectrum_away_copying}

\lean{Copying/CombinedCopyingFrontier}{41}{49}{Combined.unbounded_coherent_or_copying}

\Needspace{12\baselineskip}
\subsection{Exact application costs and host continuations}

\termdef{application-cost-handoff}
An \emph{application cost handoff} is the exact law
\[
T_n(\operatorname{before}(B,G),\mathsf{search},0)
\iff \exists r\;\bigl(n=r+6\ \land\
T_r(\operatorname{after}(B,G),\mathsf{search},0)\bigr),
\]
where $G$ satisfies the original stability premise. Apply the exact search-cost
law to the five pre-compute transitions, then application coherence to the
beta transition. Collect the six successor equations by natural-number
arithmetic. Conversely, prepend these six steps to any successful reduct run.
Existentially forgetting the exact cost gives termination equivalence between
the initial and opened-body states. In particular, failure of the reduct to
terminate cannot be hidden by application.

\lean{Copying/ApplicationCost}{9}{20}{Application.opening_cost_iff}

\lean{Copying/ApplicationCost}{24}{30}{Application.opening_terminates_iff}

For an arbitrary typed plugged host, reboundary and plug the original beta
step. The mapped source head remains application, so the step-coherence theorem
compares every possible successor with this actual plugged reduct. For an
arbitrary serial host, transport the five search steps and beta transition to
their actual host attachments. Apply the same cost decomposition to obtain
six plus the successful cost of the hosted reduct. The continuation may
compute further, diverge or get stuck; its success is not presumed.

\lean{Copying/ApplicationCost}{33}{44}{Application.host_successor}

\lean{Copying/ApplicationCost}{48}{66}{Application.host_handoff}

\Needspace{12\baselineskip}
\subsection{Regressions with unchanged translated bodies}

A translated constant with one unused parameter returns in one search step
after opening. A translated nested lambda likewise returns in one search step
without opening its inner box. Applying the six-step handoff and the exact
zero-cost law at the final state proves that both initial applications have
successful cost exactly seven, for every stable argument with the required
interface. These are all-run statements in the full combined machine.

\lean{Copying/ApplicationCoherenceExamples}{9}{17}{Application.Examples.constant_cost_iff}

\lean{Copying/ApplicationCoherenceExamples}{19}{27}{Application.Examples.nested_cost_iff}

The translated variable retains its original $D_{1,1}$ distributor. Its next
search transition after beta is a copy request, so its successful spectrum is
seven plus that request's spectrum. If the total were at most seven, the
remaining cost would be zero. A zero-step run cannot be successful at rewrite
mode, proving the no-early-success regression. This deliberately leaves the
actual remaining copying computation to the copying development.

\lean{Copying/ApplicationCoherenceExamples}{31}{38}{Application.Examples.variable_cost_iff}

\lean{Copying/ApplicationCoherenceExamples}{41}{48}{Application.Examples.variable_no_early_success}

A reference argument retains its atom attachment and satisfies the existing
stable-instance lemma, so the constant-body result specializes to successful
cost seven. Finally, the existing closed constant-body application program
has this same exact cost in the unbounded machine after every finite depth
padding: use the already proved initial-cost padding equivalence and exact
finite/unbounded successful-run correspondence, then the constant cost theorem.
Concrete divergent callbacks and repeated closure use still require the
remaining copying and template development; these results do not claim those
end-to-end regressions.

\lean{Copying/ApplicationCoherenceExamples}{51}{53}{Application.Examples.reference_argument_cost_iff}

\lean{Copying/ApplicationCoherenceExamples}{56}{60}{Application.Examples.constant_pad_cost_iff}

